\begin{table}[!htb]
\centering
\caption{Theorem~\mref{thm:bridge_envelope} case attained by each strike of every reference \RCBridge{} envelope.}
\label{tab:s_case_counts}
\footnotesize
\begin{tabular}{@{}lrrrr@{}}
\toprule
Envelope & Case (i) & Case (ii) & Case (iii) & Strikes \\
 & measured rate & assumed rate & diagnostic & \\
\midrule
E1 left-point & 0 & 7 & 0 & 7 \\
E1 \RCCE{} & 0 & 7 & 0 & 7 \\
E1 Jensen-$J$ & 1 & 4 & 2 & 7 \\
E3 reference ($H=0.10$) & 1 & 4 & 0 & 5 \\
E4 Regime~A reference & 3 & 4 & 1 & 8 \\
E4 Regime~B reference & 1 & 4 & 2 & 7 \\
E4 Regime~C reference & 2 & 5 & 0 & 7 \\
E6 (nine joint rows) & 11 & 45 & 16 & 72 \\
\midrule
Total & 19 & 80 & 21 & 120 \\
\bottomrule
\end{tabular}
\par\smallskip
\parbox{\linewidth}{\footnotesize\emph{Notes.} Case (i): the fitted rate is identifiable and used. Case (ii): the increments are submerged in noise and the envelope uses $\beta=\beta_{\mathrm{FSW}}$, which remains valid for any true rate $\beta\ge\beta_{\min}(H)$ (Section~\ref{sec:s_case_ii}). Case (iii): diagnostic display only.}
\end{table}
